\documentclass[10pt,a4paper]{amsart}
\usepackage{amsmath,amssymb,amsthm,graphicx}
% Retained preamble for the Stacks Project chapter {CH} (shared project preamble).
% The chapter file pinned-original-source.tex only carries \input{{preamble}} on its first line,
% so no preamble text lives inside the chapter itself.  The Stacks Project uses one single shared
% preamble for every chapter: the snapshot stages it three times (sources/stacks-chow/preamble.tex,
% sources/stacks-etale-cohomology/preamble.tex and focused-seeds/AGP/sources/stacks-snapshot-20261011/
% preamble.tex) and those copies are identical in size.  The block below is copied mechanically,
% line for line, from sources/stacks-chow/preamble.tex lines 14-18 and 154-217: the Xy-pic setup and
% the native theorem declarations together with the mathematical macros.  Nothing was added, removed
% or redefined; the wrapper class and the amsmath/amssymb/amsthm/graphicx base packages are supplied
% by the assembler.
% For commutative diagrams we use Xy-pic
% Chapter: stacks-cohomology.tex -- Cohomology of Sheaves.
\usepackage[all]{xy}

% We use 2cell for 2-commutative diagrams.
\xyoption{2cell}

% Theorem environments.
%
\theoremstyle{plain}
\newtheorem{theorem}[subsection]{Theorem}
\newtheorem{proposition}[subsection]{Proposition}
\newtheorem{lemma}[subsection]{Lemma}

\theoremstyle{definition}
\newtheorem{definition}[subsection]{Definition}
\newtheorem{example}[subsection]{Example}
\newtheorem{exercise}[subsection]{Exercise}
\newtheorem{situation}[subsection]{Situation}

\theoremstyle{remark}
\newtheorem{remark}[subsection]{Remark}
\newtheorem{remarks}[subsection]{Remarks}

\numberwithin{equation}{subsection}

% Macros
%
\def\lim{\mathop{\mathrm{lim}}\nolimits}
\def\colim{\mathop{\mathrm{colim}}\nolimits}
\def\Spec{\mathop{\mathrm{Spec}}}
\def\Hom{\mathop{\mathrm{Hom}}\nolimits}
\def\Ext{\mathop{\mathrm{Ext}}\nolimits}
\def\SheafHom{\mathop{\mathcal{H}\!\mathit{om}}\nolimits}
\def\SheafExt{\mathop{\mathcal{E}\!\mathit{xt}}\nolimits}
\def\Sch{\mathit{Sch}}
\def\Mor{\mathop{\mathrm{Mor}}\nolimits}
\def\Ob{\mathop{\mathrm{Ob}}\nolimits}
\def\Sh{\mathop{\mathit{Sh}}\nolimits}
\def\NL{\mathop{N\!L}\nolimits}
\def\CH{\mathop{\mathrm{CH}}\nolimits}
\def\proetale{{pro\text{-}\acute{e}tale}}
\def\etale{{\acute{e}tale}}
\def\QCoh{\mathit{QCoh}}
\def\Ker{\mathop{\mathrm{Ker}}}
\def\Im{\mathop{\mathrm{Im}}}
\def\Coker{\mathop{\mathrm{Coker}}}
\def\Coim{\mathop{\mathrm{Coim}}}

% Boxtimes
%
\DeclareMathSymbol{\boxtimes}{\mathbin}{AMSa}{"02}

%
% Macros for moduli stacks/spaces
%
\def\QCohstack{\mathcal{QC}\!\mathit{oh}}
\def\Cohstack{\mathcal{C}\!\mathit{oh}}
\def\Spacesstack{\mathcal{S}\!\mathit{paces}}
\def\Quotfunctor{\mathrm{Quot}}
\def\Hilbfunctor{\mathrm{Hilb}}
\def\Curvesstack{\mathcal{C}\!\mathit{urves}}
\def\Polarizedstack{\mathcal{P}\!\mathit{olarized}}
\def\Complexesstack{\mathcal{C}\!\mathit{omplexes}}
% \Pic is the operator that assigns to X its picard group, usage \Pic(X)
% \Picardstack_{X/B} denotes the Picard stack of X over B
% \Picardfunctor_{X/B} denotes the Picard functor of X over B
\def\Pic{\mathop{\mathrm{Pic}}\nolimits}
\def\Picardstack{\mathcal{P}\!\mathit{ic}}
\def\Picardfunctor{\mathrm{Pic}}
\def\Deformationcategory{\mathcal{D}\!\mathit{ef}}
\usepackage{cleveref}
\crefname{originallemma}{original lemma}{original lemmas}
\Crefname{originallemma}{Original lemma}{Original lemmas}
\makeatletter
\newlabel{lemma-unbounded-mayer-vietoris}{{08BX}{source}{}{sourcepointer-1}{}}
\newlabel{lemma-unbounded-mayer-vietoris@cref}{{[originallemma][0][]08BX}{[1][0][]source}{}{}{}}
\makeatother
\crefname{originallemma}{original lemma}{original lemmas}
\Crefname{originallemma}{Original lemma}{Original lemmas}
\makeatletter
\newlabel{lemma-mayer-vietoris}{{01EB}{source}{}{sourcepointer-2}{}}
\newlabel{lemma-mayer-vietoris@cref}{{[originallemma][0][]01EB}{[1][0][]source}{}{}{}}
\makeatother
\crefname{originallemma}{original lemma}{original lemmas}
\Crefname{originallemma}{Original lemma}{Original lemmas}
\makeatletter
\newlabel{lemma-factor-through-K-flat}{{0G6V}{source}{}{sourcepointer-3}{}}
\newlabel{lemma-factor-through-K-flat@cref}{{[originallemma][0][]0G6V}{[1][0][]source}{}{}{}}
\makeatother
\crefname{originallemma}{original lemma}{original lemmas}
\Crefname{originallemma}{Original lemma}{Original lemmas}
\makeatletter
\newlabel{lemma-pullback-K-flat}{{06YC}{source}{}{sourcepointer-4}{}}
\newlabel{lemma-pullback-K-flat@cref}{{[originallemma][0][]06YC}{[1][0][]source}{}{}{}}
\makeatother
\crefname{originallemma}{original lemma}{original lemmas}
\Crefname{originallemma}{Original lemma}{Original lemmas}
\makeatletter
\newlabel{more-algebra-lemma-K-flat-resolution}{{06Y4}{source}{}{sourcepointer-5}{}}
\newlabel{more-algebra-lemma-K-flat-resolution@cref}{{[originallemma][0][]06Y4}{[1][0][]source}{}{}{}}
\makeatother
\crefname{originallemma}{original lemma}{original lemmas}
\Crefname{originallemma}{Original lemma}{Original lemmas}
\makeatletter
\newlabel{modules-lemma-pullback-flat}{{05VD}{source}{}{sourcepointer-6}{}}
\newlabel{modules-lemma-pullback-flat@cref}{{[originallemma][0][]05VD}{[1][0][]source}{}{}{}}
\makeatother
\crefname{originalsection}{original section}{original sections}
\Crefname{originalsection}{Original section}{Original sections}
\makeatletter
\newlabel{derived-section-canonical-delta-functor}{{014Z}{source}{}{sourcepointer-7}{}}
\newlabel{derived-section-canonical-delta-functor@cref}{{[originalsection][0][]014Z}{[1][0][]source}{}{}{}}
\makeatother
\title{The cup product gives a map of the Mayer-Vietoris distinguished triangles}
\author{The Stacks Project Authors}
\date{stacks-project:cohomology.tex:local-snapshot-20261011; GNU FDL 1.2 or later}
\begin{document}
\maketitle
\noindent\textit{Editorial note.} Preamble: only the retained preamble of this chapter (sources/stacks-cohomology/retained-preamble.tex) is taken, because the chapter file itself carries just an input of the project's single shared preamble on its first line. It holds the Xy-pic setup, the author's native theorem declarations and the mathematical macros, copied line for line; nothing was added, removed or redefined. This card uses neither stix nor transparent, which the host TeX Live lacks. The Xy-pic matrix diagrams in the statement and in the proof rely on the xy package loaded there, which matches the two-pass local trial compilation.\par\medskip
\noindent\textit{Editorial note.} All seven external references are navigation pointers only: four same-chapter lemmas, two cross-chapter ones (the K-flat resolution in More on Algebra, the pullback flatness in Modules) and one section of Derived Categories. Their role is to resolve the original reference commands; they do not replace the construction of the four representatives and of the commutative diagram carried out inside this proof.\par\medskip
\noindent\textit{Editorial note.} Original numbering: in this chapter lemmas, equations and remarks share the subsection counter, and material-only assembly locates passages by label, so the printed numbers need not agree with the original ones.\par\medskip
\noindent\textit{Editorial note.} Uniform work/source\_title across every card of this source, as section 11 of NEW\_TASK\_AUTHORING.md requires (the manifest field work is taken from source\_title). The precise section of each claim is carried by claim\_key and by statement\_anchor in the registration, and the section name is repeated in the editorial notes where it matters.\par\medskip
\noindent\textit{Allowed standard prerequisites (editorial).}\begin{itemize}
\item The unbounded Mayer-Vietoris distinguished triangles for a union of two opens. Same chapter, Lemma labelled lemma-unbounded-mayer-vietoris (canonical Stacks tag 08BX, source line 8200), named in the counted statement at lines 8349 and 8353.. Supplies the two triangles that the cup-product maps are required to compare; the comparison itself is the counted content.
\item Existence of K-flat and K-injective representatives with flat, respectively injective, terms. Same chapter: Lemma lemma-pullback-K-flat (tag 06YC, line 5961) and Lemma lemma-factor-through-K-flat (tag 0G6V, line 6233); cross-chapter: More on Algebra, Lemma more-algebra-lemma-K-flat-resolution (tag 06Y4).. Lets the proof replace K and M by complexes on which the total tensor product is honestly computed, so that the comparison maps can be written on the level of complexes.
\item Functoriality of the distinguished triangle associated to a short exact sequence of complexes of A-modules. Derived Categories, section labelled derived-section-canonical-delta-functor (canonical Stacks tag 014Z), used at proof lines 8428-8430.. Converts the commutative diagram of short exact sequences of complexes into the required map between the two distinguished triangles.
\end{itemize}\medskip
\ifdefined\hypertarget\hypertarget{sourcepointer-1}{}\fi
\noindent\textit{Source reference.} Original lemma 08BX: sources/stacks-cohomology/pinned-original-source.tex, source lines 8200--8200. Same-chapter pointer to the distinguished triangle of the unbounded Mayer-Vietoris sequence, named three times in the counted statement (lines 8349 and 8353) and used again at proof line 8425. It supplies the two triangles whose functoriality is compared; the comparison maps themselves (choice of T\textasciicircum{}bullet, K\textasciicircum{}bullet, I\textasciicircum{}bullet, J\textasciicircum{}bullet and the commutative diagram of complexes of A-modules) are constructed inside the extracted proof.\par\medskip
\ifdefined\hypertarget\hypertarget{sourcepointer-2}{}\fi
\noindent\textit{Source reference.} Original lemma 01EB: sources/stacks-cohomology/pinned-original-source.tex, source lines 767--767. Same-chapter pointer invoked at proof lines 8420-8424 only for the fact that the columns of the displayed diagram are exact sequences of complexes of A-modules. That exactness is the standard two-open Mayer-Vietoris input; the counted content of this card is the comparison of cup products, not the Mayer-Vietoris sequence.\par\medskip
\ifdefined\hypertarget\hypertarget{sourcepointer-3}{}\fi
\noindent\textit{Source reference.} Original lemma 0G6V: sources/stacks-cohomology/pinned-original-source.tex, source lines 6233--6233. Same-chapter pointer applied at proof lines 8369-8379 to replace the adjunction map T\textasciicircum{}bullet \textbackslash{}otimes\_A O\_X -> K by a map of complexes into a K-flat complex with flat terms. A standard K-flat replacement step; the surrounding construction stays in the extracted proof.\par\medskip
\ifdefined\hypertarget\hypertarget{sourcepointer-4}{}\fi
\noindent\textit{Source reference.} Original lemma 06YC: sources/stacks-cohomology/pinned-original-source.tex, source lines 5961--5961. Same-chapter pointer used at proof line 8367 for flatness of the pullback T\textasciicircum{}bullet \textbackslash{}otimes\_A O\_X, a standard prerequisite for the bookkeeping of flat terms.\par\medskip
\ifdefined\hypertarget\hypertarget{sourcepointer-5}{}\fi
\noindent\textit{Source reference.} Original lemma 06Y4: sources/stacks-chow/tags.txt, source lines 8873--8873. Cross-chapter pointer (More on Algebra) used at proof line 8360 to choose a K-flat complex of flat A-modules representing R\textbackslash{}Gamma(X, K); a standard resolution prerequisite, not part of the counted comparison argument.\par\medskip
\ifdefined\hypertarget\hypertarget{sourcepointer-6}{}\fi
\noindent\textit{Source reference.} Original lemma 05VD: sources/stacks-chow/tags.txt, source lines 7547--7547. Cross-chapter pointer (Modules) used at proof line 8368 for flatness of the pullback of a flat module; standard prerequisite.\par\medskip
\ifdefined\hypertarget\hypertarget{sourcepointer-7}{}\fi
\noindent\textit{Source reference.} Original section 014Z: sources/stacks-chow/tags.txt, source lines 1490--1490. Cross-chapter navigation pointer (Derived Categories, section on the canonical delta functor) used at proof lines 8428-8430 to pass from a short exact sequence of complexes to its associated distinguished triangle. Only that functoriality is quoted; the triangles and the comparison maps are built in the extracted proof.\par\medskip
\begin{lemma}
\label{lemma-mayer-vietoris-cup}
Let $(X, \mathcal{O}_X)$ be a ringed space. Set $A = \Gamma(X, \mathcal{O}_X)$.
Suppose that $X = U \cup V$ is a union of two open subsets. For objects
$K$ and $M$ of $D(\mathcal{O}_X)$ we have a map of distinguished triangles
$$
\xymatrix{
R\Gamma(X, K) \otimes_A^\mathbf{L} R\Gamma(X, M) \ar[r] \ar[d] &
R\Gamma(X, K \otimes_{\mathcal{O}_X}^\mathbf{L} M) \ar[d] \\
R\Gamma(X, K) \otimes_A^\mathbf{L}
(R\Gamma(U, M) \oplus R\Gamma(V, M)) \ar[r] \ar[d] &
R\Gamma(U, K \otimes_{\mathcal{O}_X}^\mathbf{L} M)
\oplus R\Gamma(V, K \otimes_{\mathcal{O}_X}^\mathbf{L} M)) \ar[d] \\
R\Gamma(X, K) \otimes_A^\mathbf{L} R\Gamma(U \cap V, M) \ar[r] \ar[d] &
R\Gamma(U \cap V, K \otimes_{\mathcal{O}_X}^\mathbf{L} M) \ar[d] \\
R\Gamma(X, K) \otimes_A^\mathbf{L} R\Gamma(X, M)[1] \ar[r] &
R\Gamma(X, K \otimes_{\mathcal{O}_X}^\mathbf{L} M)[1]
}
$$
where
\begin{enumerate}
\item the horizontal arrows are given by cup product,
\item on the right hand side we have the distinguished triangle
of Lemma \ref{lemma-unbounded-mayer-vietoris} for
$K \otimes_{\mathcal{O}_X}^\mathbf{L} M$, and
\item on the left hand side we have the exact functor
$R\Gamma(X, K) \otimes_A^\mathbf{L} - $ applied to the
distinguished triangle of Lemma \ref{lemma-unbounded-mayer-vietoris} for $M$.
\end{enumerate}
\end{lemma}

\begin{proof}
Choose a K-flat complex $T^\bullet$ of flat $A$-modules representing
$R\Gamma(X, K)$, see
More on Algebra, Lemma \ref{more-algebra-lemma-K-flat-resolution}.
Denote $T^\bullet \otimes_A \mathcal{O}_X$ the pullback of $T^\bullet$
by the morphism of ringed spaces $(X, \mathcal{O}_X) \to (pt, A)$.
There is a natural adjunction map
$\epsilon : T^\bullet \otimes_A \mathcal{O}_X \to K$ in $D(\mathcal{O}_X)$.
Observe that $T^\bullet \otimes_A \mathcal{O}_X$ is a K-flat
complex of $\mathcal{O}_X$-modules with flat terms, see
Lemma \ref{lemma-pullback-K-flat} and
Modules, Lemma \ref{modules-lemma-pullback-flat}.
By Lemma \ref{lemma-factor-through-K-flat} we can find a morphism of complexes
$$
T^\bullet \otimes_A \mathcal{O}_X \longrightarrow \mathcal{K}^\bullet
$$
of $\mathcal{O}_X$-modules representing $\epsilon$
such that $\mathcal{K}^\bullet$ is a
K-flat complex with flat terms. Namely, by the construction of
$D(\mathcal{O}_X)$ we can first represent $\epsilon$ by some map of complexes
$e : T^\bullet \otimes_A \mathcal{O}_X \to \mathcal{L}^\bullet$
of $\mathcal{O}_X$-modules representing $\epsilon$
and then we can apply the lemma to $e$. Choose a K-injective
complex $\mathcal{I}^\bullet$ whose terms are injective $\mathcal{O}_X$-modules
representing $M$. Finally, choose a quasi-isomorphism
$$
\text{Tot}(\mathcal{K}^\bullet \otimes_\mathcal{O} \mathcal{I}^\bullet)
\longrightarrow
\mathcal{J}^\bullet
$$
into a K-injective complex whose terms are injective $\mathcal{O}_X$-modules.
Observe that source and target of this arrow represent
$K \otimes_{\mathcal{O}_X}^\mathbf{L} M$ in $D(\mathcal{O}_X)$.
At this point, for any open $W \subset X$ we obtain a map of complexes
$$
\text{Tot}(T^\bullet \otimes_A \mathcal{I}^\bullet(W))
\to
\text{Tot}(\mathcal{K}^\bullet(W) \otimes_A \mathcal{I}^\bullet(W))
\to
\mathcal{J}^\bullet(W)
$$
of $A$-modules whose composition represents the map
$$
R\Gamma(X, K) \otimes_A^\mathbf{L} R\Gamma(W, M)
\longrightarrow
R\Gamma(W, K \otimes_{\mathcal{O}_X}^\mathbf{L} M)
$$
in $D(A)$. Clearly, these maps are compatible with restriction mappings.
OK, so now we can consider the following commutative(!) diagram
of complexes of $A$-modules
$$
\xymatrix{
0 \ar[d] & 0 \ar[d] \\
\text{Tot}(T^\bullet \otimes_A \mathcal{I}^\bullet(X)) \ar[d] \ar[r] &
\mathcal{J}^\bullet(X) \ar[d] \\
\text{Tot}(T^\bullet \otimes_A
(\mathcal{I}^\bullet(U) \oplus \mathcal{I}^\bullet(V)) \ar[d] \ar[r] &
\mathcal{J}^\bullet(U) \oplus \mathcal{J}^\bullet(V) \ar[d] \\
\text{Tot}(T^\bullet \otimes_A \mathcal{I}^\bullet(U \cap V)) \ar[r] \ar[d] &
\mathcal{J}^\bullet(U \cap V) \ar[d] \\
0 & 0
}
$$
By the proof of Lemma \ref{lemma-mayer-vietoris} the columns are
exact sequences of complexes of $A$-modules (this also uses that
$\text{Tot}(T^\bullet \otimes_A -)$ transforms short exact sequences
of complexes of $A$-modules into short exact sequences as the terms
of $T^\bullet$ are flat $A$-modules). Since the distinguished triangles
of Lemma \ref{lemma-unbounded-mayer-vietoris}
are the distinguished triangles associated to these
short exact sequences of complexes, the desired result follows from
the functoriality of ``taking the associated distinguished triangle''
discussed in
Derived Categories, Section \ref{derived-section-canonical-delta-functor}.
\end{proof}

\end{document}
